\section{Definitions}
\label{section-definition-categories}

\noindent
We recall the definitions, partly to fix notation.

\begin{definition}
\label{definition-category}
A {\it category} $\mathcal{C}$ consists of the following data:
\begin{enumerate}
\item A set of objects $\Ob(\mathcal{C})$.
\item For each pair $x, y \in \Ob(\mathcal{C})$ a set of morphisms
$\Mor_\mathcal{C}(x, y)$.
\item For each triple $x, y, z\in \Ob(\mathcal{C})$ a composition
map $ \Mor_\mathcal{C}(y, z) \times \Mor_\mathcal{C}(x, y)
\to \Mor_\mathcal{C}(x, z) $, denoted $(\phi, \psi) \mapsto
\phi \circ \psi$.
\end{enumerate}
These data are to satisfy the following rules:
\begin{enumerate}
\item For every element $x\in \Ob(\mathcal{C})$ there exists a
morphism $\text{id}_x\in \Mor_\mathcal{C}(x, x)$ such that
$\text{id}_x \circ \phi = \phi$ and $\psi \circ \text{id}_x = \psi $ whenever
these compositions make sense.
\item Composition is associative, i.e., $(\phi \circ \psi) \circ \chi =
\phi \circ ( \psi \circ \chi)$ whenever these compositions make sense.
\end{enumerate}
\end{definition}

\noindent
It is customary to require all the morphism sets
$\Mor_\mathcal{C}(x, y)$ to be disjoint.
In this way a morphism $\phi : x \to y$ has a unique {\it source} $x$
and a unique {\it target} $y$. This is not strictly necessary,
although care has to be taken in formulating condition (2) above
if it is not the case. It is convenient and we will often assume
this is the case. In this case we say that $\phi$ and $\psi$ are
{\it composable} if the source of $\phi$ is equal to the
target of $\psi$, in which case $\phi \circ \psi$ is defined.
An equivalent definition would be to define a category
as a quintuple $(\text{Ob}, \text{Arrows}, s, t, \circ)$
consisting of a set of objects, a set of morphisms (arrows),
source, target and composition subject to a long list of axioms.
We will occasionally use this point of view.

\begin{remark}
\label{remark-big-categories}
Big categories. In some texts a category is allowed to have a proper
class of objects. We will allow this as well in these notes but only
in the following list of cases (to be updated as we go along).
In particular, when we say: ``Let $\mathcal{C}$ be a category''
then it is understood that $\Ob(\mathcal{C})$ is a set.
\begin{enumerate}
\item The category $\textit{Sets}$ of sets.
\item The category $\textit{Ab}$ of abelian groups.
\item The category $\textit{Groups}$ of groups.
\item Given a group $G$ the category $G\textit{-Sets}$ of
sets with a left $G$-action.
\item Given a ring $R$ the category $\text{Mod}_R$ of $R$-modules.
\item Given a field $k$ the category of vector spaces over $k$.
\item The category of rings.
\item The category of divided power rings, see
Divided Power Algebra, Section \ref{dpa-section-divided-power-rings}.
\item The category of schemes.
\item The category $\textit{Top}$ of topological spaces.
\item Given a topological space $X$ the category
$\textit{PSh}(X)$ of presheaves of sets over $X$.
\item Given a topological space $X$ the category
$\Sh(X)$ of sheaves of sets over $X$.
\item Given a topological space $X$ the category
$\textit{PAb}(X)$ of presheaves of abelian groups over $X$.
\item Given a topological space $X$ the category
$\textit{Ab}(X)$ of sheaves of abelian groups over $X$.
\item Given a small category $\mathcal{C}$ the category of functors
from $\mathcal{C}$ to $\textit{Sets}$.
\item Given a category $\mathcal{C}$ the category of presheaves of sets
over $\mathcal{C}$.
\item Given a site $\mathcal{C}$ the category of sheaves
of sets over $\mathcal{C}$.
\end{enumerate}
One of the reason to enumerate these here is to try and avoid
working with something like the ``collection'' of ``big'' categories
which would be like working with the collection of all classes
which I think definitively is a meta-mathematical object.
\end{remark}

\begin{remark}
\label{remark-unique-identity}
It follows directly from the definition that any two identity morphisms
of an object $x$ of $\mathcal{A}$ are the same. Thus we may and will
speak of {\it the} identity morphism $\text{id}_x$ of $x$.
\end{remark}

\begin{definition}
\label{definition-isomorphism}
A morphism $\phi : x \to y$ is an {\it isomorphism} of the category
$\mathcal{C}$ if there exists a morphism $\psi : y \to x$
such that $\phi \circ \psi = \text{id}_y$ and
$\psi \circ \phi = \text{id}_x$.
\end{definition}

\noindent
An isomorphism $\phi$ is also sometimes called an {\it invertible}
morphism, and the morphism $\psi$ of the definition is called the
{\it inverse} and denoted $\phi^{-1}$. It is unique if it exists. Note that
given an object $x$ of a category $\mathcal{A}$ the set of invertible
elements $\text{Aut}_\mathcal{A}(x)$
of $\Mor_\mathcal{A}(x, x)$ forms a group under composition.
This group is called the {\it automorphism} group of $x$ in $\mathcal{A}$.

\begin{definition}
\label{definition-groupoid}
A {\it groupoid} is a category where every morphism is an isomorphism.
\end{definition}

\begin{example}
\label{example-group-groupoid}
A group $G$ gives rise to a groupoid with a single object $x$
and morphisms $\Mor(x, x) = G$, with the composition rule
given by the group law in $G$. Every groupoid with a single
object is of this form.
\end{example}

\begin{example}
\label{example-set-groupoid}
A set $C$ gives rise to a groupoid $\mathcal{C}$ defined as follows:
As objects we take $\Ob(\mathcal{C}) := C$ and for morphisms
we take $\Mor(x, y)$ empty if $x\neq y$ and equal to
$\{\text{id}_x\}$ if $x = y$.
\end{example}

\begin{definition}
\label{definition-functor}
A {\it functor} $F : \mathcal{A} \to \mathcal{B}$
between two categories $\mathcal{A}, \mathcal{B}$ is given by the
following data:
\begin{enumerate}
\item A map $F : \Ob(\mathcal{A}) \to \Ob(\mathcal{B})$.
\item For every $x, y \in \Ob(\mathcal{A})$ a map
$F : \Mor_\mathcal{A}(x, y) \to \Mor_\mathcal{B}(F(x), F(y))$,
denoted $\phi \mapsto F(\phi)$.
\end{enumerate}
These data should be compatible with composition and identity morphisms
in the following manner: $F(\phi \circ \psi) =
F(\phi) \circ F(\psi)$ for a composable pair $(\phi, \psi)$ of
morphisms of $\mathcal{A}$ and $F(\text{id}_x) = \text{id}_{F(x)}$.
\end{definition}

\noindent
Note that every category $\mathcal{A}$ has an
{\it identity} functor $\text{id}_\mathcal{A}$.
In addition, given a functor $G : \mathcal{B} \to \mathcal{C}$
and a functor $F : \mathcal{A} \to \mathcal{B}$ there is
a {\it composition} functor $G \circ F : \mathcal{A} \to \mathcal{C}$
defined in an obvious manner.

\begin{definition}
\label{definition-faithful}
Let $F : \mathcal{A} \to \mathcal{B}$ be a functor.
\begin{enumerate}
\item We say $F$ is {\it faithful} if
for any objects $x, y \in \Ob(\mathcal{A})$ the map
$$
F : \Mor_\mathcal{A}(x, y) \to \Mor_\mathcal{B}(F(x), F(y))
$$
is injective.
\item If these maps are all bijective then $F$ is called
{\it fully faithful}.
\item
The functor $F$ is called {\it essentially surjective} if for any
object $y \in \Ob(\mathcal{B})$ there exists an object
$x \in \Ob(\mathcal{A})$ such that $F(x)$ is isomorphic to $y$ in
$\mathcal{B}$.
\end{enumerate}
\end{definition}

\begin{definition}
\label{definition-subcategory}
A {\it subcategory} of a category $\mathcal{B}$ is a category $\mathcal{A}$
whose objects and arrows form subsets of the objects and arrows of $\mathcal{B}$
and such that source, target and composition in $\mathcal{A}$ agree with those
of $\mathcal{B}$ and such that the identity morphism of an object of
$\mathcal{A}$ matches the one in $\mathcal{B}$. We say $\mathcal{A}$ is a
{\it full subcategory} of $\mathcal{B}$ if $\Mor_\mathcal{A}(x, y)
= \Mor_\mathcal{B}(x, y)$ for all $x, y \in \Ob(\mathcal{A})$.
We say $\mathcal{A}$ is a {\it strictly full} subcategory of $\mathcal{B}$
if it is a full subcategory and given $x \in \Ob(\mathcal{A})$ any
object of $\mathcal{B}$ which is isomorphic to $x$ is also in $\mathcal{A}$.
\end{definition}

\noindent
If $\mathcal{A} \subset \mathcal{B}$ is a subcategory then the
identity map is a functor from $\mathcal{A}$ to $\mathcal{B}$.
Furthermore a subcategory $\mathcal{A} \subset \mathcal{B}$
is full if and only if the inclusion functor is fully faithful.
Note that given a category $\mathcal{B}$ the set of full subcategories
of $\mathcal{B}$ is the same as the set of subsets of
$\Ob(\mathcal{B})$.

\begin{remark}
\label{remark-functor-into-sets}
Suppose that $\mathcal{A}$ is a category.
A functor $F$ from $\mathcal{A}$ to $\textit{Sets}$
is a mathematical object (i.e., it is a set not a class or a formula
of set theory, see
Sets, Section \ref{sets-section-sets-everything})
even though the category of sets is ``big''.
Namely, the range of $F$ on objects will be
a set $F(\Ob(\mathcal{A}))$ and then we
may think of $F$ as a functor between
$\mathcal{A}$ and the full subcategory
of the category of sets whose
objects are elements of $F(\Ob(\mathcal{A}))$.
\end{remark}

\begin{example}
\label{example-group-homomorphism-functor}
A homomorphism $p : G\to H$ of groups gives rise to a functor
between the associated groupoids in Example \ref{example-group-groupoid}. It is
faithful (resp.\ fully faithful) if and only if $p$ is injective (resp.\ an
isomorphism).
\end{example}

\begin{example}
\label{example-category-over-X}
Given a category $\mathcal{C}$ and an object $X\in \Ob(\mathcal{C})$
we define the {\it category of objects over $X$},
denoted $\mathcal{C}/X$ as follows.
The objects of $\mathcal{C}/X$ are morphisms $Y\to X$ for
some $Y\in \Ob(\mathcal{C})$. Morphisms between objects
$Y\to X$ and $Y'\to X$ are morphisms $Y\to Y'$ in $\mathcal{C}$ that
make the obvious diagram commute. Note that there is a functor
$p_X : \mathcal{C}/X\to \mathcal{C}$ which simply forgets the
morphism. Moreover given a morphism $f : X'\to X$ in
$\mathcal{C}$ there is an induced functor
$F : \mathcal{C}/X' \to \mathcal{C}/X$ obtained by composition with $f$,
and $p_X\circ F = p_{X'}$.
\end{example}

\begin{example}
\label{example-category-under-X}
Given a category $\mathcal{C}$ and an object $X\in \Ob(\mathcal{C})$
we define the {\it category of objects under $X$},
denoted $X/\mathcal{C}$ as follows.
The objects of $X/\mathcal{C}$ are morphisms $X\to Y$ for
some $Y\in \Ob(\mathcal{C})$. Morphisms between objects
$X\to Y$ and $X\to Y'$ are morphisms $Y\to Y'$ in $\mathcal{C}$ that
make the obvious diagram commute. Note that there is a functor
$p_X : X/\mathcal{C}\to \mathcal{C}$ which simply forgets the
morphism. Moreover given a morphism $f : X'\to X$ in
$\mathcal{C}$ there is an induced functor
$F : X/\mathcal{C} \to X'/\mathcal{C}$
obtained by composition with $f$,
and $p_{X'}\circ F = p_X$.
\end{example}




\begin{definition}
\label{definition-transformation-functors}
Let $F, G : \mathcal{A} \to \mathcal{B}$ be functors.
A {\it natural transformation}, or a {\it morphism of functors}
$t : F \to G$, is a collection $\{t_x\}_{x\in \Ob(\mathcal{A})}$
such that
\begin{enumerate}
\item $t_x : F(x) \to G(x)$ is a morphism in the category $\mathcal{B}$, and
\item for every morphism $\phi : x \to y$ of $\mathcal{A}$ the following
diagram is commutative
$$
\xymatrix{
F(x) \ar[r]^{t_x} \ar[d]_{F(\phi)} & G(x) \ar[d]^{G(\phi)} \\
F(y) \ar[r]^{t_y} & G(y) }
$$
\end{enumerate}
\end{definition}

\noindent
Sometimes we use the diagram
$$
\xymatrix{
\mathcal{A}
\rtwocell^F_G{t}
&
\mathcal{B}
}
$$
to indicate that $t$ is a morphism from $F$ to $G$.

\medskip\noindent
Note that every functor $F$ comes with the {\it identity} transformation
$\text{id}_F : F \to F$. In addition, given a morphism of
functors $t : F \to G$ and a morphism of functors $s : E \to F$
then the {\it composition} $t \circ s$ is defined by the rule
$$
(t \circ s)_x = t_x \circ s_x : E(x) \to G(x)
$$
for $x \in \Ob(\mathcal{A})$.
It is easy to verify that this is indeed a morphism of functors
from $E$ to $G$.
In this way, given categories
$\mathcal{A}$ and $\mathcal{B}$ we obtain a new category,
namely the category of functors between $\mathcal{A}$ and
$\mathcal{B}$.

\begin{remark}
\label{remark-functors-sets-sets}
This is one instance where the same thing does not hold if
$\mathcal{A}$ is a ``big'' category. For example consider
functors $\textit{Sets} \to \textit{Sets}$. As we have currently
defined it such a functor is a class and not a set. In other
words, it is given by a formula in set theory (with some variables
equal to specified sets)! It is not a good idea to try to consider
all possible formulae of set theory as part of the definition of
a mathematical object. The same problem presents itself when
considering sheaves on the category of schemes for example.
We will come back to this point later.
\end{remark}

\begin{definition}
\label{definition-equivalence-categories}
An {\it equivalence of categories}
$F : \mathcal{A} \to \mathcal{B}$ is a functor such that there
exists a functor $G : \mathcal{B} \to \mathcal{A}$ such that
the compositions $F \circ G$ and $G \circ F$ are isomorphic to the
identity functors $\text{id}_\mathcal{B}$,
respectively $\text{id}_\mathcal{A}$.
In this case we say that $G$ is a {\it quasi-inverse} to $F$.
\end{definition}

\begin{lemma}
\label{lemma-construct-quasi-inverse}
Let $F : \mathcal{A} \to \mathcal{B}$ be a fully faithful functor.
Suppose for every $X \in \Ob(\mathcal{B})$ we are given an
object $j(X)$ of $\mathcal{A}$ and an isomorphism $i_X : X \to F(j(X))$.
Then there is a unique functor $j : \mathcal{B} \to \mathcal{A}$
such that $j$ extends the rule on objects, and the isomorphisms
$i_X$ define an isomorphism of functors
$\text{id}_\mathcal{B} \to F \circ j$. Moreover, $j$ and $F$
are quasi-inverse equivalences of categories.
\end{lemma}

\begin{proof}
To construct $j : \mathcal{B} \to \mathcal{A}$, there are two steps. Firstly,
we define the map $j : \Ob(\mathcal{B}) \to \Ob(\mathcal{A})$ that associates
$j(X)$ to $X \in \mathcal{B}$. Secondly, if $X,Y \in \Ob(\mathcal{B})$ and
$\phi : X \to Y$, we consider $\phi' := i_Y \circ\phi\circ i_X^{-1}$. There is
a unique $\varphi$ verifying $F(\varphi) = \phi'$, using that $F$ is fully
faithful. We define $j(\phi) = \varphi$. We omit the verification that $j$
is a functor. By construction the diagram
$$
\xymatrix{
X \ar[r]_-{i_X} \ar[d]_{\phi}
&
F(j(X)) \ar[d]^{F\circ j(\phi)}
\\
Y \ar[r]^-{i_Y}
&
F(j(Y))
}
$$
commutes. Hence, as each $i_X$ is an isomorphism, $\{i_X\}_X$ is an
isomorphism of functors $\text{id}_\mathcal{B} \to F\circ j$. To conclude, we
have to also prove that $j \circ F$ is isomorphic to $\text{id}_\mathcal{A}$.
However, since $F$ is fully faithful, in order to do this it suffices
to prove this after post-composing with $F$, i.e., it suffices to
show that $F \circ j \circ F$ is isomorphic to $F \circ \text{id}_\mathcal{A}$
(small detail omitted). Since $F \circ j \cong \text{id}_\mathcal{B}$
this is clear.
\end{proof}

\begin{lemma}
\label{lemma-equivalence-categories}
A functor is an equivalence of categories if and only if it is both fully
faithful and essentially surjective.
\end{lemma}

\begin{proof}
Let $F : \mathcal{A} \to \mathcal{B}$ be essentially surjective and fully
faithful. As by convention all categories are small and as $F$ is essentially
surjective we can, using the axiom of choice, choose for every
$X \in \Ob(\mathcal{B})$ an object $j(X)$ of $\mathcal{A}$ and an
isomorphism $i_X : X \to F(j(X))$. Then we apply
Lemma \ref{lemma-construct-quasi-inverse}
using that $F$ is fully faithful.
\end{proof}

\begin{definition}
\label{definition-product-category}
Let $\mathcal{A}$, $\mathcal{B}$ be categories.
We define the {\it product category}
$\mathcal{A} \times \mathcal{B}$ to be the category with
objects
$\Ob(\mathcal{A} \times \mathcal{B}) =
\Ob(\mathcal{A}) \times \Ob(\mathcal{B})$
and
$$
\Mor_{\mathcal{A} \times \mathcal{B}}((x, y), (x', y'))
:=
\Mor_\mathcal{A}(x, x')\times
\Mor_\mathcal{B}(y, y').
$$
Composition is defined componentwise.
\end{definition}


